\begin{afterword}
\summaryhead

The post takes the author's earlier treatment of the falling-tree riddle as its model. By the author's account, that treatment did not answer the riddle. Instead it explained the mental process that makes it feel like a question.

The post describes a series of philosophers. Some answer questions such as ``Do we have free will?'', some define their terms first, and some argue that the question is meaningless. None of them, it says, explains why people have the intuition in the first place. The author suggests that the feeling of an open question corresponds to a central unit in a neural network that stays unsettled after every observable fact is fixed.

The essay then separates ``explaining why'' a belief exists, as an evolutionary story does, from ``explaining how'' the mind produces it. Only the second dissolves a question. You are done when no confusion remains. You can feel done without being done, but real dissolution is unmistakable. The post ends by setting free will as homework: describe the cognitive algorithm that produces the debate.

\responsehead

The post places its author above philosophy. It lines up philosophers, each wiser than the last, and puts its own method at the top. No philosopher is named.

Its one substantial idea is that an intuition, even a mistaken one, is a fact about minds that needs explaining. The idea is real, and it is old. Wittgenstein wrote that the clarity he aimed at ``simply means that the philosophical problems should completely disappear'' (\textsc{Philosophical Investigations}, §133). For free will in particular, Daniel Wegner had published a cognitive theory of how the feeling of willing arises in 2002. The post presents dissolution as a path its author found alone, ``once upon a time.''

The standard it sets is severe: a ``stack trace,'' ``in detail,'' ``step by step.'' It does not meet that standard itself. Its mechanism, the ``dangling unit,'' is qualified in a parenthesis as ``metaphorically speaking.'' It gives no evidence that brains contain any such thing. The example it points to for ``explaining how'' is a diagram in another post. By its own test it is a just-so story about mechanism. The post condemns exactly that kind of story when the story is evolutionary.

The test of success is a feeling. The post argues that feelings come from algorithms that need not match the world. It then trusts one feeling, the feeling of being done. It calls that feeling ``unmistakable,'' and grants in the next line that one can feel done without being done. A reader who follows this post has no way to tell an explanation that dissolved their confusion from one that merely satisfied them. The post warns against that very satisfaction, but only when it comes from refutations.

The homework forbids arguing ``that free will is compatible with determinism, or not.'' Three months later the author's own solution appeared, and it contained such a thesis. Agency and moral responsibility, it said, ``require determinism.'' Hobart had argued this in \textsc{Mind} in 1934. The post puts that answer out of bounds for its readers, and its author then gave it under a new name.

In short: a demand for step-by-step explanation, made in a post whose own mechanism is an admitted metaphor, which scores success by a feeling it admits can mislead, and which forbids its readers the position its author went on to take.
\end{afterword}
